\begin{abstract}

Zero-knowledge virtual machines (zkVMs) prove correct execution of RISC-V programs, but their security depends on faithful alignment between the RISC-V specification and the arithmetization that certifies an execution trace. When this alignment breaks, a prover can forge proofs for incorrect executions, undermining the trust model of systems securing billions of dollars. Detecting such \emph{semantic gaps} is challenging: each zkVM encodes its trace differently, and bugs cluster at rare boundary conditions that are effectively unreachable by random testing.

Our key insight is that soundness bugs in zkVMs fall into two structurally distinct classes---executor \emph{divergences} and constraint \emph{underconstraints}---and that, despite heterogeneous internal representations, the boundary conditions triggering both classes are shared across backends. We present \tool, a semantics-guided differential fuzzing framework that exploits this structure. \tool defines a \emph{cross-zkVM semantic model} that lifts six production backends into a common space of 17 observation types and 32 semantic categories, and builds a \emph{boundary-driven fuzzing framework} where category signature novelty replaces edge coverage, a multi-armed bandit governs eight ISA-aware mutation strategies, and semantic witness injection probes underconstrained circuits at identified boundary points. We establish a two-phase coverage characterization showing that the detection phases are complementary: neither subsumes the other, and removing witness injection reduces recall from 81.8\% to at most 11.4\%.

Evaluated across OpenVM, SP1, RISC Zero, Jolt, Nexus, and Pico on 44 known vulnerabilities, \tool achieves 81.8\% recall (36/44). On the four backends where the prior state-of-the-art (Arguzz) has full support, \tool detects 1.6$\times$ more bugs (18 vs.\ 11); across all six backends, the improvement is 3.3$\times$ (36 vs.\ 11). \tool additionally discovers 6 previously unknown zero-day vulnerabilities including critical soundness violations in production systems. All findings have been responsibly disclosed and patched.

\end{abstract}
